% C26 idade por serie
\begin{table}[htbp]\centering
\caption{Idade m\'edia (ponderada) por s\'erie e ano, alunos do EM regular (V3003A=06) frequentando no $Q_2$. Teste cr\'itico: se a repet\^encia real tivesse aumentado em 2022--2023, a idade m\'edia por s\'erie deveria ter aumentado.}
\label{tab:idade_serie}
\begin{threeparttable}\small
\begin{tabular}{lrrrrrr}\toprule
Ano & 1\textsuperscript{o} EM & 2\textsuperscript{o} EM & 3\textsuperscript{o} EM & 5\textsuperscript{o} EF & 9\textsuperscript{o} EF & N$^{a}$ EM\\
\midrule
2017 & 15.9 & 16.9 & 18.1 & 10.7 & 14.6 & 25,059 \\
2018 & 15.9 & 16.8 & 18.1 & 10.5 & 14.6 & 23,096 \\
2019 & 15.9 & 16.8 & 18.0 & 10.5 & 14.6 & 22,882 \\
2020$^{\dagger}$ & 15.9 & 16.8 & 18.3 & 10.5 & 14.5 & 15,431 \\
2021$^{\dagger}$ & 16.1 & 17.0 & 19.2 & 11.0 & 14.7 & 15,561 \\
2022 & 15.9 & 16.8 & 18.6 & 10.7 & 14.6 & 20,159 \\
2023 & 15.8 & 16.7 & 18.3 & 10.5 & 14.7 & 19,057 \\
2024 & 15.6 & 16.6 & 17.8 & 10.4 & 14.5 & 19,042 \\
\bottomrule\end{tabular}
\begin{tablenotes}\footnotesize
\item \textit{Fonte:} PNADC trimestral, $Q_2$ de cada ano. Idade m\'edia ponderada por $V1028$.
\item $^{a}$N total para 1\textsuperscript{o} + 2\textsuperscript{o} + 3\textsuperscript{o} EM em $Q_2$.
\item $^{\dagger}$Anos COVID-19.
\end{tablenotes}\end{threeparttable}\end{table}
